\chapter{第3章：Hecke算子 习题}

\difficultykey

\section{基础题 \easy}

\begin{problem}{计算四次级数的像}
\easy
设
\[
E_4(\tau)=1+240q+2160q^2+6720q^3+17520q^4+\cdots.
\]
利用权 $4$ 的 Hecke 算子在 $q$-展开上的公式，计算 $T_2(E_4)$ 的常数项以及 $q,q^2,q^3$ 的系数。
\end{problem}
\begin{solution}
对权 $4$ 模形式，
\[
T_2f=\sum_{n\ge0}\bigl(a_{2n}+2^3a_{n/2}\bigr)q^n.
\]
这里 $a_{n/2}=0$，若 $n/2\notin \mathbf Z_{\ge0}$。因此
\[
[q^0]T_2(E_4)=a_0+8a_0=9.
\]
又有
\[
[q]T_2(E_4)=a_2=2160,
\]
\[
[q^2]T_2(E_4)=a_4+8a_1=17520+8\cdot240=19440,
\]
\[
[q^3]T_2(E_4)=a_6.
\]
而 $a_6=240\sigma_3(6)=240(1+2^3+3^3+6^3)=240\cdot252=60480$。故
\[
T_2(E_4)=9+2160q+19440q^2+60480q^3+\cdots.
\]
事实上这正好等于 $9E_4$ 的前几项，因为 $E_4$ 是 Hecke 特征形式。
\end{solution}

\begin{problem}{读取首项特征值}
\easy
设 $f(\tau)=\sum_{n\ge1}a_nq^n\in S_k(\Gamma)$，并满足 $T_pf=\lambda_pf$。证明 $q$ 的系数满足 $a_p=\lambda_pa_1$。特别地，若 $f$ 归一化为 $a_1=1$，则 $a_p=\lambda_p$。
\end{problem}
\begin{solution}
由
\[
T_pf=\sum_{n\ge1}\bigl(a_{pn}+p^{k-1}a_{n/p}\bigr)q^n
\]
可知 $q$ 的系数来自 $n=1$，此时第二项消失，所以
\[
[q]T_pf=a_p.
\]
另一方面，$\lambda_pf$ 的 $q$ 系数是 $\lambda_pa_1$。比较两边的 $q$ 系数即得
\[
a_p=\lambda_pa_1.
\]
若 $a_1=1$，便得到 $a_p=\lambda_p$。
\end{solution}

\begin{problem}{计算判别形式的像}
\easy
设
\[
\Delta(\tau)=q-24q^2+252q^3-1472q^4+4830q^5-6048q^6+\cdots.
\]
利用 Hecke 公式计算 $T_2(\Delta)$ 的前三个非零项。
\end{problem}
\begin{solution}
$\Delta$ 的权为 $12$，故
\[
T_2\Delta=\sum_{n\ge1}\bigl(\tau(2n)+2^{11}\tau(n/2)\bigr)q^n.
\]
于是
\[
[q]T_2\Delta=\tau(2)=-24,
\]
\[
[q^2]T_2\Delta=\tau(4)+2^{11}\tau(1)=-1472+2048=576,
\]
\[
[q^3]T_2\Delta=\tau(6)=-6048.
\]
所以
\[
T_2\Delta=-24q+576q^2-6048q^3+\cdots.
\]
因为 $\Delta$ 是归一化特征形式，这也等于 $\tau(2)\Delta$ 的前几项。
\end{solution}

\begin{problem}{检验互素乘法}
\easy
利用 $\Delta$ 的前几项，验证在这一个例子上
\[
T_2T_3\Delta=T_6\Delta.
\]
只需比较 $q,q^2$ 两项即可。
\end{problem}
\begin{solution}
先算
\[
T_3\Delta=252\Delta=252q-6048q^2+63504q^3+\cdots.
\]
再对它施以 $T_2$，利用权 $12$ 的公式，得到
\[
[q]T_2T_3\Delta=[q^2]T_3\Delta=-6048,
\]
\[
[q^2]T_2T_3\Delta=[q^4]T_3\Delta+2^{11}[q]T_3\Delta.
\]
由于 $[q^4]T_3\Delta=252\tau(4)=252\cdot(-1472)=-370944$，且 $[q]T_3\Delta=252$，所以
\[
[q^2]T_2T_3\Delta=-370944+2048\cdot252=145152.
\]
另一方面，$T_6\Delta=\tau(6)\Delta=-6048\Delta$，故
\[
[q]T_6\Delta=-6048,
\qquad
[q^2]T_6\Delta=-6048\cdot(-24)=145152.
\]
前两项一致，符合互素情形的乘法公式 $T_2T_3=T_6$。
\end{solution}

\begin{problem}{验证平方递推}
\easy
已知 $\tau(2)=-24$，$\tau(4)=-1472$。验证归一化特征形式在素数平方处的关系
\[
\tau(4)=\tau(2)^2-2^{11}.
\]
\end{problem}
\begin{solution}
直接代入即可：
\[
\tau(2)^2-2^{11}=(-24)^2-2048=576-2048=-1472=	au(4).
\]
这正是一般关系
\[
a_{p^2}=a_p^2-p^{k-1}
\]
在 $\Delta\in S_{12}(\mathrm{SL}_2(\mathbf Z))$ 上的具体实例。
\end{solution}

\begin{problem}{说明菱形算子良定义}
\easy
设 $f\in M_k(\Gamma_1(N))$，取 $d\in (\mathbf Z/N\mathbf Z)^\times$。若
\[
\langle d\rangle f=f\bigm|_k\gamma,
\qquad
\gamma=\begin{pmatrix}a&b\\Nc&d'\end{pmatrix}\in\Gamma_0(N),\quad d'\equiv d\pmod N,
\]
证明此定义与 $\gamma$ 的选取无关。
\end{problem}
\begin{solution}
设 $\gamma_1,\gamma_2\in\Gamma_0(N)$ 的右下角都与 $d$ 同余模 $N$。则
\[
\delta:=\gamma_2\gamma_1^{-1}\in\Gamma_1(N).
\]
原因是：$\delta$ 的左下角仍被 $N$ 整除，而对角元都与 $1$ 同余模 $N$。于是
\[
f\bigm|_k\gamma_2=f\bigm|_k\delta\gamma_1=(f\bigm|_k\delta)\bigm|_k\gamma_1=f\bigm|_k\gamma_1,
\]
因为 $f$ 对 $\Gamma_1(N)$ 不变。故 $\langle d\rangle f$ 良定义。
\end{solution}

\begin{problem}{验证菱形算子的乘法}
\easy
证明对任意 $d_1,d_2\in (\mathbf Z/N\mathbf Z)^\times$，都有
\[
\langle d_1\rangle\langle d_2\rangle=\langle d_1d_2\rangle.
\]
\end{problem}
\begin{solution}
取 $\gamma_i\in\Gamma_0(N)$，使其右下角分别同余于 $d_i$。则
\[
\langle d_i\rangle f=f\bigm|_k\gamma_i.
\]
于是
\[
\langle d_1\rangle\langle d_2\rangle f=(f\bigm|_k\gamma_2)\bigm|_k\gamma_1=f\bigm|_k(\gamma_2\gamma_1).
\]
而 $\gamma_2\gamma_1\in\Gamma_0(N)$ 的右下角模 $N$ 同余于 $d_2d_1=d_1d_2$，故按定义
\[
f\bigm|_k(\gamma_2\gamma_1)=\langle d_1d_2\rangle f.
\]
因此结论成立。特别地，菱形算子给出一个阿贝尔群表示。
\end{solution}

\begin{problem}{保持尖点条件}
\easy
说明若 $f\in S_k(\mathrm{SL}_2(\mathbf Z))$，则 $T_pf$ 仍是尖点形式。
\end{problem}
\begin{solution}
尖点形式在无穷尖点的 $q$-展开没有常数项，即 $a_0=0$。而
\[
T_pf=\sum_{n\ge0}\bigl(a_{pn}+p^{k-1}a_{n/p}\bigr)q^n,
\]
所以它的常数项是
\[
a_0+p^{k-1}a_0=0.
\]
故 $T_pf$ 在无穷尖点仍消没。对 $\mathrm{SL}_2(\mathbf Z)$ 只有这一个尖点，因此 $T_pf$ 仍为尖点形式，即
\[
T_p\bigl(S_k(\mathrm{SL}_2(\mathbf Z))\bigr)\subseteq S_k(\mathrm{SL}_2(\mathbf Z)).
\]
\end{solution}

\begin{problem}{读出局部因子}
\easy
设 $f=\sum_{n\ge1}a_nq^n$ 是权 $k$ 的归一化 Hecke 特征形式。写出 $L(f,s)$ 在素数 $p$ 处的 Euler 因子，并解释它只由 $a_p$ 与权 $k$ 决定。
\end{problem}
\begin{solution}
对归一化 Hecke 特征形式，章节中给出的 Euler 乘积为
\[
L(f,s)=\prod_p\left(1-a_pp^{-s}+p^{k-1-2s}\right)^{-1}.
\]
因此 $p$ 处的局部因子就是
\[
L_p(f,s)=\left(1-a_pp^{-s}+p^{k-1-2s}\right)^{-1}.
\]
可见它只依赖于该素数处的 Hecke 特征值 $a_p$，以及权带来的固定项 $p^{k-1}$。一旦知道所有 $a_p$，整条 Euler 乘积便被确定。
\end{solution}

\section{进阶题 \medium}

\begin{problem}{写出一般系数公式}
\medium
设 $f(\tau)=\sum_{n\ge0}a_nq^n\in M_k(\mathrm{SL}_2(\mathbf Z))$。证明 $T_mf$ 的 $q^n$ 系数可写成
\[
\sum_{d\mid (m,n)} d^{k-1}a_{mn/d^2}.
\]
\end{problem}
\begin{solution}
先对素数幂 $m=p^r$ 说明。由反复应用
\[
T_pf=\sum_{n\ge0}(a_{pn}+p^{k-1}a_{n/p})q^n
\]
以及递推关系
\[
T_{p^{r+1}}=T_pT_{p^r}-p^{k-1}T_{p^{r-1}},
\]
可归纳得到
\[
[q^n]T_{p^r}f=\sum_{j=0}^{\min(r,v_p(n))}p^{j(k-1)}a_{p^{r+n_p-2j}n'}
\]
其中 $n=p^{n_p}n'$，$(n',p)=1$。这正是
\[
\sum_{d\mid (p^r,n)} d^{k-1}a_{p^rn/d^2}
\]
的写法。

再把一般的 $m$ 分解为互素素数幂乘积，利用互素时 $T_{ab}=T_aT_b$，即可把各素数贡献相乘，得到总公式
\[
[q^n]T_mf=\sum_{d\mid (m,n)} d^{k-1}a_{mn/d^2}.
\]
\end{solution}

\begin{problem}{展开素数幂算子}
\medium
利用递推式
\[
T_{p^{r+1}}=T_pT_{p^r}-p^{k-1}T_{p^{r-1}}
\]
把 $T_{p^2}$ 与 $T_{p^3}$ 都写成 $T_p$ 的多项式。
\end{problem}
\begin{solution}
取 $r=1$，立得
\[
T_{p^2}=T_p^2-p^{k-1}.
\]
再取 $r=2$，有
\[
T_{p^3}=T_pT_{p^2}-p^{k-1}T_p.
\]
代入上一式，得到
\[
T_{p^3}=T_p(T_p^2-p^{k-1})-p^{k-1}T_p=T_p^3-2p^{k-1}T_p.
\]
所以前两项分别是
\[
T_{p^2}=T_p^2-p^{k-1},
\qquad
T_{p^3}=T_p^3-2p^{k-1}T_p.
\]
\end{solution}

\begin{problem}{推出互素乘性}
\medium
设 $f=\sum_{n\ge1}a_nq^n$ 是归一化 Hecke 特征形式。证明当 $(m,n)=1$ 时，
\[
a_{mn}=a_ma_n.
\]
\end{problem}
\begin{solution}
因为 $f$ 是归一化特征形式，$T_rf=a_rf$ 对每个 $r\ge1$ 都成立。若 $(m,n)=1$，则
\[
T_mT_nf=T_{mn}f.
\]
左边作用在 $f$ 上给出
\[
T_mT_nf=T_m(a_nf)=a_nT_mf=a_na_mf.
\]
右边则是
\[
T_{mn}f=a_{mn}f.
\]
于是
\[
a_{mn}f=a_ma_nf.
\]
比较 $q$ 的系数，或者直接利用 $f\neq0$，就得到
\[
a_{mn}=a_ma_n.
\]
\end{solution}

\begin{problem}{推出素数幂递推}
\medium
设 $f=\sum_{n\ge1}a_nq^n$ 是权 $k$ 的归一化 Hecke 特征形式。证明对任意素数 $p$ 与 $r\ge1$，
\[
a_{p^{r+1}}=a_pa_{p^r}-p^{k-1}a_{p^{r-1}}.
\]
\end{problem}
\begin{solution}
由 Hecke 算子递推，
\[
T_{p^{r+1}}=T_pT_{p^r}-p^{k-1}T_{p^{r-1}}.
\]
把它作用到 $f$ 上，得到
\[
a_{p^{r+1}}f=a_pa_{p^r}f-p^{k-1}a_{p^{r-1}}f.
\]
因为 $f\neq0$，可消去 $f$，于是
\[
a_{p^{r+1}}=a_pa_{p^r}-p^{k-1}a_{p^{r-1}}.
\]
这就是素数幂系数的标准递推。
\end{solution}

\begin{problem}{计算若干拉马努金值}
\medium
利用归一化特征形式关系，计算 $\tau(8)$ 与 $\tau(9)$。
\end{problem}
\begin{solution}
对 $\Delta$ 有权 $k=12$，且
\[
\tau(4)=\tau(2)^2-2^{11}=-1472.
\]
再用一次递推，
\[
\tau(8)=\tau(2)\tau(4)-2^{11}\tau(2).
\]
代入 $\tau(2)=-24$、$\tau(4)=-1472$，得
\[
\tau(8)=(-24)(-1472)-2048(-24)=35328+49152=84480.
\]
同理，
\[
\tau(9)=\tau(3)^2-3^{11}.
\]
因为 $\tau(3)=252$，所以
\[
\tau(9)=252^2-177147=63504-177147=-113643.
\]
故答案为
\[
\tau(8)=84480,
\qquad
\tau(9)=-113643.
\]
\end{solution}

\begin{problem}{利用自伴性得正交性}
\medium
设 $f,g\in S_k(\Gamma)$ 都是 $T_n$ 的特征向量，且特征值分别为 $\lambda_n,\mu_n$。若 $\lambda_n\ne\mu_n$，证明 $\langle f,g\rangle=0$。
\end{problem}
\begin{solution}
由 Petersson 内积的自伴性，
\[
\langle T_nf,g\rangle=\langle f,T_ng\rangle.
\]
把特征关系代入，得
\[
\lambda_n\langle f,g\rangle=\mu_n\langle f,g\rangle.
\]
移项可得
\[
(\lambda_n-\mu_n)\langle f,g\rangle=0.
\]
若 $\lambda_n\ne\mu_n$，则只能有
\[
\langle f,g\rangle=0.
\]
因此，不同特征值的特征向量彼此正交。
\end{solution}

\begin{problem}{证明正交补稳定}
\medium
设 $V\subset S_k(\Gamma)$ 是 Hecke 算子稳定子空间。证明它的 Petersson 正交补 $V^{\perp}$ 也对所有 Hecke 算子稳定。
\end{problem}
\begin{solution}
取 $h\in V^{\perp}$，要证 $T_nh\in V^{\perp}$。任取 $v\in V$，由自伴性
\[
\langle T_nh,v\rangle=\langle h,T_nv\rangle.
\]
因为 $V$ 对 $T_n$ 稳定，所以 $T_nv\in V$。又 $h\in V^{\perp}$，于是
\[
\langle h,T_nv\rangle=0.
\]
故对任意 $v\in V$ 都有 $\langle T_nh,v\rangle=0$，这说明
\[
T_nh\in V^{\perp}.
\]
因此 $V^{\perp}$ 也是 Hecke 稳定的。
\end{solution}

\begin{problem}{构造最简单旧形式}
\medium
设 $f\in S_k(\Gamma_0(N))$，且素数 $p\nmid N$。证明 $f(\tau)$ 与 $f(p\tau)$ 都属于 $S_k(\Gamma_0(Np))$，并解释它们生成了从水平 $N$ 升到水平 $Np$ 的最简单旧形式子空间。
\end{problem}
\begin{solution}
因为 $\Gamma_0(Np)\subseteq\Gamma_0(N)$，故任何 $\Gamma_0(N)$ 上的尖点形式自动也是 $\Gamma_0(Np)$ 上的尖点形式，所以 $f(\tau)\in S_k(\Gamma_0(Np))$。

再看 $f(p\tau)$。它的 $q$-展开是
\[
f(p\tau)=\sum_{n\ge1}a_nq^{pn},
\]
因此在无穷尖点仍无常数项，是尖点形式。章节中的退化映射理论说明，$f(\tau)$ 与 $f(p\tau)$ 正是从水平 $N$ 到水平 $Np$ 的两条标准嵌入，因此它们张成旧形式空间中的一个二维候选子空间，通常记为由 $f$ 产生的 oldclass。
\end{solution}

\begin{problem}{形式推导欧拉乘积}
\medium
设 $f=\sum_{n\ge1}a_nq^n$ 是归一化 Hecke 特征形式，并满足互素乘性与素数幂递推。形式地证明
\[
\sum_{n\ge1}\frac{a_n}{n^s}=\prod_p\left(1-a_pp^{-s}+p^{k-1-2s}\right)^{-1}.
\]
\end{problem}
\begin{solution}
先固定一个素数 $p$，考虑局部级数
\[
\sum_{r\ge0}a_{p^r}X^r.
\]
由递推
\[
a_{p^{r+1}}=a_pa_{p^r}-p^{k-1}a_{p^{r-1}}
\]
可知它满足
\[
\left(1-a_pX+p^{k-1}X^2\right)\sum_{r\ge0}a_{p^r}X^r=1.
\]
故
\[
\sum_{r\ge0}a_{p^r}X^r=\frac{1}{1-a_pX+p^{k-1}X^2}.
\]
取 $X=p^{-s}$，得到
\[
\sum_{r\ge0}\frac{a_{p^r}}{p^{rs}}=\left(1-a_pp^{-s}+p^{k-1-2s}\right)^{-1}.
\]
再利用互素乘性，任意正整数的系数可按素因子分解，故整体 Dirichlet 级数分解为各素数局部因子的乘积，即
\[
L(f,s)=\prod_p\left(1-a_pp^{-s}+p^{k-1-2s}\right)^{-1}.
\]
\end{solution}

\begin{problem}{讨论绝对收敛区域}
\medium
设 $f=\sum_{n\ge1}a_nq^n$ 满足估计 $|a_n|\le Cn^A$。证明其 Dirichlet 级数
\[
L(f,s)=\sum_{n\ge1}\frac{a_n}{n^s}
\]
在 $\Re(s)>A+1$ 绝对收敛。
\end{problem}
\begin{solution}
记 $s=\sigma+it$。则
\[
\left|\frac{a_n}{n^s}\right|=\frac{|a_n|}{n^\sigma}\le Cn^{A-\sigma}.
\]
若 $\sigma>A+1$，则指数 $A-\sigma<-1$，从而比较级数
\[
\sum_{n\ge1}n^{A-\sigma}
\]
收敛。于是
\[
\sum_{n\ge1}\left|\frac{a_n}{n^s}\right|
\]
收敛，也就是说 $L(f,s)$ 在半平面 $\Re(s)>A+1$ 绝对收敛。
\end{solution}

\begin{problem}{证明Hecke代数交换}
\medium
利用互素乘法公式与素数幂递推，证明一切 Hecke 算子彼此交换。
\end{problem}
\begin{solution}
先看同一个素数 $p$ 的幂。由
\[
T_{p^2}=T_p^2-p^{k-1},
\qquad
T_{p^3}=T_p^3-2p^{k-1}T_p,
\]
以及一般递推
\[
T_{p^{r+1}}=T_pT_{p^r}-p^{k-1}T_{p^{r-1}},
\]
可归纳得每个 $T_{p^r}$ 都是 $T_p$ 的整系数多项式，所以它们彼此交换。

若 $(m,n)=1$，则
\[
T_mT_n=T_{mn}=T_nT_m.
\]
把一般的 $m,n$ 写成素数幂分解后，不同素数部分两两互素，因此也彼此交换。综合起来便得
\[
T_mT_n=T_nT_m
\qquad (m,n\ge1).
\]
\end{solution}

\begin{problem}{由特征值恢复系数}
\medium
设 $f=\sum_{n\ge1}a_nq^n$ 是归一化 Hecke 特征形式。证明只要知道所有素数处的特征值 $a_p$，就能递推出全部系数 $a_n$。
\end{problem}
\begin{solution}
因为 $f$ 归一化，$a_1=1$。若知道所有素数处的 $a_p$，则由素数幂递推
\[
a_{p^{r+1}}=a_pa_{p^r}-p^{k-1}a_{p^{r-1}}
\]
可以依次求出所有 $a_{p^r}$。接着，任意正整数可写作
\[
n=\prod_i p_i^{r_i}.
\]
利用互素乘性，得
\[
a_n=\prod_i a_{p_i^{r_i}}.
\]
因此每个 $a_n$ 都由各个素数处的 $a_p$ 唯一决定。
\end{solution}

\section{挑战题 \hard}

\begin{problem}{同时对角化Hecke作用}
\hard
设 $V\subset S_k(\Gamma)$ 是有限维、且对所有 Hecke 算子稳定的复向量子空间。证明可以在 $V$ 中找到一组 Petersson 正交基，使得每个基向量都是所有 $T_n$ 的共同特征向量。
\end{problem}
\begin{solution}
由章节结论，每个 $T_n$ 都是 Petersson 内积下的自伴算子，因此在有限维厄米空间 $V$ 上可正交对角化。

先对 $T_1=\mathrm{id}$ 以外任选一个非平凡 Hecke 算子，例如 $T_2$。把 $V$ 分解成 $T_2$ 的互相正交特征子空间之和。由于所有 Hecke 算子彼此交换，每个 $T_n$ 都保持这些特征子空间稳定。于是我们只需在每个特征子空间上继续对角化 $T_3$。重复此过程，就得到越来越细的正交分解。

因为 $V$ 有限维，只需处理有限个步骤便可得到由共同特征子空间组成的直和分解。每个一维共同特征子空间取一个单位向量，就得到所求的 Petersson 正交基。故 $V$ 可由共同 Hecke 特征形式构成正交基。
\end{solution}

\begin{problem}{证明特征值为实数}
\hard
设 $f\in S_k(\Gamma)$ 是某个 Hecke 算子 $T_n$ 的非零特征向量，满足 $T_nf=\lambda f$。证明 $\lambda\in\mathbf R$。
\end{problem}
\begin{solution}
由自伴性，
\[
\langle T_nf,f\rangle=\langle f,T_nf\rangle.
\]
把特征关系代入左边与右边，得到
\[
\lambda\langle f,f\rangle=\overline{\lambda}\langle f,f\rangle.
\]
这里 $\langle f,f\rangle>0$，因为 Petersson 内积正定且 $f\neq0$。故
\[
\lambda=\overline{\lambda}.
\]
所以 $\lambda$ 是实数。
\end{solution}

\begin{problem}{写出局部母函数}
\hard
设 $f=\sum_{n\ge1}a_nq^n$ 是权 $k$ 的归一化 Hecke 特征形式。证明对任意素数 $p$，都有
\[
\sum_{r\ge0}a_{p^r}X^r=\frac{1}{1-a_pX+p^{k-1}X^2}.
\]
\end{problem}
\begin{solution}
记
\[
A_p(X)=\sum_{r\ge0}a_{p^r}X^r.
\]
其中 $a_{p^0}=a_1=1$。由递推关系
\[
a_{p^{r+1}}=a_pa_{p^r}-p^{k-1}a_{p^{r-1}}\qquad (r\ge1)
\]
可得
\[
A_p(X)-1-a_pX=\sum_{r\ge1}a_{p^{r+1}}X^{r+1}
\]
\[
=a_pX\sum_{r\ge1}a_{p^r}X^r-p^{k-1}X^2\sum_{r\ge1}a_{p^{r-1}}X^{r-1}
\]
\[
=a_pX(A_p(X)-1)-p^{k-1}X^2A_p(X).
\]
整理即得
\[
A_p(X)\bigl(1-a_pX+p^{k-1}X^2\bigr)=1.
\]
故
\[
A_p(X)=\frac{1}{1-a_pX+p^{k-1}X^2}.
\]
\end{solution}

\begin{problem}{由素数数据决定形式}
\hard
设 $f=\sum_{n\ge1}a_nq^n$ 与 $g=\sum_{n\ge1}b_nq^n$ 是同一空间中的两个归一化 Hecke 特征形式，并且对每个素数 $p$ 都有 $a_p=b_p$。证明 $f=g$。
\end{problem}
\begin{solution}
因为两者都归一化，$a_1=b_1=1$。由上一章和本章的 Hecke 理论，素数幂系数满足相同的递推
\[
a_{p^{r+1}}=a_pa_{p^r}-p^{k-1}a_{p^{r-1}},
\qquad
b_{p^{r+1}}=b_pb_{p^r}-p^{k-1}b_{p^{r-1}}.
\]
又已知 $a_p=b_p$，故归纳可得对每个 $r\ge0$ 都有
\[
a_{p^r}=b_{p^r}.
\]
任意整数 $n=\prod p_i^{r_i}$，再由互素乘性得到
\[
a_n=\prod_i a_{p_i^{r_i}}=\prod_i b_{p_i^{r_i}}=b_n.
\]
因此 $f$ 与 $g$ 的全部 Fourier 系数相同，所以
\[
f=g.
\]
这就是该章多重性一定理的最基本形式。
\end{solution}

\begin{problem}{总结拉马努金递推}
\hard
证明 Ramanujan 函数满足
\[
\tau(mn)=\tau(m)\tau(n)\quad ((m,n)=1),
\]
以及
\[
\tau(p^{r+1})=\tau(p)\tau(p^r)-p^{11}\tau(p^{r-1}).
\]
并由此计算 $\tau(12)$ 与 $\tau(18)$。
\end{problem}
\begin{solution}
$\Delta$ 是权 $12$ 的归一化 Hecke 特征形式，因此一般理论直接给出
\[
\tau(mn)=\tau(m)\tau(n)\quad ((m,n)=1),
\]
和
\[
\tau(p^{r+1})=\tau(p)\tau(p^r)-p^{11}\tau(p^{r-1}).
\]
对 $12=3\cdot4$，因为 $(3,4)=1$，故
\[
\tau(12)=\tau(3)\tau(4)=252\cdot(-1472)=-370944.
\]
对 $18=2\cdot9$，因为 $(2,9)=1$，故
\[
\tau(18)=\tau(2)\tau(9)=-24\cdot(-113643)=2727432.
\]
这说明所有 $\tau(n)$ 都可由素数处数据逐步恢复。
\end{solution}

\begin{problem}{定义新空间}
\hard
设 $S_k(\Gamma_0(N))_{\mathrm{old}}$ 是由较低水平通过退化映射得到的旧形式空间。说明如何把新空间定义为其 Petersson 正交补，并证明这个新空间对所有 Hecke 算子稳定。
\end{problem}
\begin{solution}
按照 Atkin--Lehner 理论，先把所有真因子水平 $M\mid N$ 的尖点形式，通过标准退化映射嵌入到 $S_k(\Gamma_0(N))$ 中；这些像张成旧形式空间
\[
S_k(\Gamma_0(N))_{\mathrm{old}}.
\]
然后定义新空间为其 Petersson 正交补：
\[
S_k(\Gamma_0(N))_{\mathrm{new}}:=S_k(\Gamma_0(N))_{\mathrm{old}}^{\perp}.
\]
由于旧形式空间本身由 Hecke 稳定子空间张成，因此它对所有 Hecke 算子稳定。上一节已经证明：任一 Hecke 稳定子空间的正交补仍是 Hecke 稳定的。所以
\[
T_n\bigl(S_k(\Gamma_0(N))_{\mathrm{new}}\bigr)\subseteq S_k(\Gamma_0(N))_{\mathrm{new}}
\]
对每个 $n\ge1$ 都成立。故新空间继承了 Hecke 作用。
\end{solution}

\begin{problem}{分析一维空间的作用}
\hard
设某个尖点形式空间 $S_k(\Gamma)$ 是一维的，且 $f$ 是其中的归一化非零元素。证明对每个 $n\ge1$，都有
\[
T_nf=a_nf.
\]
\end{problem}
\begin{solution}
因为 $S_k(\Gamma)$ 一维，所以任意线性算子在这空间上都只能按某个标量作用。特别地，对每个 $n$ 都存在常数 $\lambda_n$，使得
\[
T_nf=\lambda_nf.
\]
再比较 $q$ 的系数。左边的 $q$ 系数是 $a_n$，右边的 $q$ 系数是 $\lambda_na_1$。由于 $f$ 归一化，$a_1=1$，所以
\[
\lambda_n=a_n.
\]
于是
\[
T_nf=a_nf.
\]
这说明在一维尖点空间里，归一化生成元自动就是 Hecke 特征形式。
\end{solution}

\begin{problem}{构造正交特征基}
\hard
设 $V\subset S_k(\Gamma)$ 是有限维 Hecke 稳定子空间。证明 $V$ 可以分解成若干两两正交的共同特征子空间之直和。
\end{problem}
\begin{solution}
由 Hecke 算子交换且都自伴，上一题已经说明 $V$ 中可以找到共同特征向量构成的正交基。把具有完全相同特征值系统的向量收集在一起，就得到共同特征子空间。

不同特征系统的向量至少在某个 $T_n$ 上具有不同特征值，因此由正交性判别可知这些子空间彼此正交。它们张成整个 $V$，故
\[
V=\bigoplus_{\chi} V_{\chi}
\]
其中 $V_{\chi}$ 表示给定 Hecke 特征值系统 $\chi$ 的共同特征子空间。于是 $V$ 分解为若干两两正交的共同特征子空间之直和。
\end{solution}
